\begin{lemma}\label{lmm:precompact}
Let $(\mc Q, d)$ be a complete metric space and $A\subset \mc Q$. Assume that all closed subsets $B\subset A$ are compact. Then $\overline A$ is compact.
\end{lemma}
\begin{proof}
	As $(\mc Q, d)$ is complete and $\overline A$ is closed, $\overline A$ is complete. Thus, it suffices to show that $\overline A$ is totally bounded, which is equivalent to $A$ being totally bounded.

	Assume $A$ is not totally bounded. Then there is $\epsilon > 0$ such that no finite collection of balls of radius $\epsilon$ covers $A$. Hence, we can recursively choose $a_1 \in A$ and $a_{k+1} \in A \setminus \bigcup_{j=1}^{k}\ball_\epsilon(a_j)$ for $k \in \N$. The resulting set $S = \Set{a_k \given k \in\N}$ is infinite and $\epsilon$-separated, i.e., $d(a_j, a_k) \geq \epsilon$ for all $j \neq k$. In particular, every ball of radius $\frac12\epsilon$ contains at most one element of $S$. Thus, $S$ has no point of accumulation in $\mc Q$, which makes $S$ a closed subset of $\mc Q$, and the cover of $S$ by the balls $\ball_{\frac12\epsilon}(a_k)$, $k\in\N$, has no finite subcover, which makes $S$ non-compact. As $S \subset A$, this contradicts the assumption of the lemma.

	Hence, $A$ is totally bounded and $\overline A$ is compact.
\end{proof}
